\documentclass[a4paper,12pt]{article}
\usepackage[utf8]{inputenc}
\usepackage[russian]{babel}
\usepackage{amsmath}
\usepackage{graphicx}
\usepackage{geometry}
\geometry{left=2cm, right=2cm, top=2cm, bottom=2cm}
\usepackage{caption}
\usepackage{hyperref}
\usepackage{setspace}
\usepackage{indentfirst}
\usepackage{titlesec}

% GOST-like formatting
\onehalfspacing          % 1.5 line spacing
\setlength{\parindent}{1.25cm} % paragraph indent
\titleformat{\section}[block]{\bfseries\large}{}{0pt}{}
\titleformat{\subsection}[block]{\bfseries\normalsize}{}{0pt}{}

\begin{document}

\begin{center}
    {\Large \textbf{Раздел 1: НЕПРЕРЫВНО-ДЕТЕРМИНИРОВАННЫЙ ПОДХОД}}
\end{center}

\vspace{10mm}
% Title
\begin{center}
    \textbf{\large Часть 1 -- Модель с постоянными интенсивностями}
\end{center}
\vspace{5mm}

\section*{Часть 1. Непрерывно-детерминированная модель с постоянными интенсивностями}

\subsection*{Цель}
Смоделировать передачу битового потока от источника к приёмнику с постоянными скоростями:
поступление $\lambda = 17$ бит/с, обработка $\mu = 10$ бит/с. Буфер неограничен, начальное
состояние $q(0)=0$.

\subsection*{Параметры модели}
Основные параметры модели сведены в таблицу~\ref{tab:params}.

\begin{table}[h!]
    \centering
    \caption{Параметры модели}
    \label{tab:params}
    \begin{tabular}{|l|c|l|}
        \hline
        Параметр & Обозначение & Значение \\
        \hline
        Интенсивность поступления & $\lambda$ & 17 бит/с \\
        Интенсивность обработки   & $\mu$    & 10 бит/с \\
        Максимальный размер буфера & $b$     & $\infty$ (не ограничен) \\
        Начальная очередь         & $q(0)$  & 0 бит \\
        Время моделирования       & $T$     & 10 с \\
        \hline
    \end{tabular}
\end{table}

\subsection*{Математическая модель}
Динамика очереди описывается дифференциальным уравнением первого порядка
\[
\frac{dq(t)}{dt} = \lambda - \mu = 7\ \text{бит/с}, \qquad q(0)=0,
\]
откуда аналитическое решение $q(t) = 7t$.

\subsection*{Реализация в Simulink}
Основные компоненты и их
настройки:
\begin{enumerate}
    \item \textbf{Constant (Lambda)} – значение $17$.
    \item \textbf{Constant (Mu)} – значение $10$.
    \item \textbf{Sum (dq\_dt)} – вычисляет разность $\lambda - \mu$ (входы $+-$).
    \item \textbf{Integrator (Buffer)} – интегрирует разность с начальным условием $0$.
          Включено ограничение выхода: нижний предел $0$ (очередь не может быть отрицательной),
          верхний предел $\infty$ (буфер неограничен).
    \item \textbf{Clock} и \textbf{Gain} (коэффициент $7$) – формируют теоретическую прямую
          $7t$ для сравнения с результатом моделирования.
    \item \textbf{Scope (Scope\_Queue)} – двухвходовой осциллограф; на первый вход подаётся
          $q(t)$, на второй – теоретическая прямая $7t$.
    \item \textbf{Display} и \textbf{To Workspace} – для наблюдения текущего значения
          и экспорта данных.
\end{enumerate}

\subsection*{Схема модели}
На рис.~\ref{fig:sheme_start} представлена блок-схема модели в начальный момент времени ($t = 0$).
На рис.~\ref{fig:sheme_end} – та же схема после завершения моделирования ($t = 10$ с),
когда в буфере накопилось 70 бит.

\begin{figure}[!ht]
    \centering
    \includegraphics[width=\textwidth]{scheme_start.png}   % <-- скриншот схемы при t = 0
    \caption{Схема модели в начальный момент времени ($t = 0$ с).}
    \label{fig:sheme_start}
\end{figure}

\begin{figure}[h!]
    \centering
    \includegraphics[width=\textwidth]{scheme_end.png}     % <-- скриншот схемы при t = 10 с
    \caption{Схема модели после 10 секунд моделирования ($t = 10$ с).}
    \label{fig:sheme_end}
\end{figure}

\subsection*{График состояния буфера}
На рис.~\ref{fig:graph} приведён график изменения длины очереди во времени, полученный с
помощью блока Scope.

\begin{figure}[h!]
    \centering
    \includegraphics[width=0.85\textwidth]{graph.png}      % <-- скриншот графика Scope
    \caption{График зависимости $q(t)$ от времени.}
    \label{fig:graph}
\end{figure}

\subsection*{Анализ}
Из уравнения $\frac{dq}{dt} = 7$ бит/с следует линейный рост очереди с наклоном $7$.
При $t = 10$ с теоретическое значение $q = 7 \times 10 = 70$ бит. Сравнение
рис.~\ref{fig:sheme_start} и \ref{fig:sheme_end} подтверждает, что за 10 секунд длина очереди
выросла от $0$ до $70$ бит, что полностью совпадает с аналитическим решением. Ограничение
интегратора снизу нулём гарантирует, что очередь никогда не принимает отрицательных значений.
Так как верхнее ограничение отключено, буфер ведёт себя как бесконечный – рост очереди не
прекращается.

\subsection*{Вывод по части 1}
При постоянных интенсивностях $\lambda > \mu$ и неограниченном буфере очередь неограниченно
возрастает по линейному закону $q(t) = (\lambda - \mu)\,t$. Результаты моделирования полностью
соответствуют аналитическому решению, что подтверждает корректность непрерывно-
детерминированного подхода для данной задачи.

\newpage\section*{Часть 2. Модель с постоянными интенсивностями и тайм-аутом (дискретно-событийный анализ)}

\subsection*{Цель}
В систему из Части~1 вводится ограничение времени ожидания каждого бита в очереди:
\[
t_{\text{lim}} = \frac{17}{10^{2}} = 0.17\ \text{с}.
\]
Если бит ждёт начала обслуживания дольше $t_{\text{lim}}$, он безвозвратно отбрасывается.  
Требуется провести аналитический расчёт первых отбрасываний, проследить динамику очереди
на уровне отдельных заявок и получить графики состояния буфера и накопленного числа
отброшенных битов.

\subsection*{Параметры модели}
\begin{table}[h!]
    \centering
    \caption{Параметры модели с тайм-аутом (Часть~2)}
    \label{tab:params2}
    \begin{tabular}{|l|c|l|}
        \hline
        Параметр & Обозначение & Значение \\
        \hline
        Интенсивность поступления & $\lambda$ & 17 бит/с \\
        Интенсивность обработки   & $\mu$    & 10 бит/с \\
        Длительность обслуживания & $1/\mu$  & 0.1 с \\
        Допустимое время ожидания  & $t_{\text{lim}}$ & 0.17 с \\
        Начальная очередь         & $q(0)$  & 0 бит \\
        Время моделирования       & $T$     & 15 с \\
        \hline
    \end{tabular}
\end{table}

\subsection*{Математическая модель}
В отличие от жидкостного приближения, здесь каждый бит рассматривается как отдельная
заявка. Процесс функционирования системы описывается следующими правилами:
\begin{itemize}
    \item Биты поступают строго периодически: момент поступления $n$-го бита
          $t_n^{\text{arr}} = n / \lambda = n/17$~с.
    \item Обслуживание одного бита занимает ровно $1/\mu = 0.1$~с.
    \item Буфер организован по дисциплине FIFO. Если в момент освобождения сервера
          очередь не пуста, на обслуживание становится первый ожидающий бит.
    \item Для каждого бита $i$, находящегося в очереди, непрерывно контролируется
          время ожидания: если $t - t_i^{\text{arr}} > t_{\text{lim}}$, бит
          немедленно удаляется из буфера (сгорает) и в дальнейшем не обслуживается.
\end{itemize}

\subsubsection*{Аналитический расчёт первых битов}
Проведём пошаговую трассировку событий для начального отрезка времени.

\begin{table}[h!]
    \centering
    \caption{Трассировка первых битов}
    \label{tab:bit_trace}
    \begin{tabular}{|c|c|c|c|c|c|}
        \hline
        \textbf{Бит} & \textbf{Прибытие, с} & \textbf{Начало обслуж., с} & \textbf{Конец обслуж., с} & \textbf{Ожидание, с} & \textbf{Статус} \\
        \hline
        0 & 0.0000 & 0.0000 & 0.1000 & 0.0000 & обслужен \\
        1 & 0.0588 & 0.1000 & 0.2000 & 0.0412 & обслужен \\
        2 & 0.1176 & 0.2000 & 0.3000 & 0.0824 & обслужен \\
        3 & 0.1765 & 0.3000 & 0.4000 & 0.1235 & обслужен \\
        4 & 0.2353 & 0.4000 & 0.5000 & 0.1647 & обслужен \\
        5 & 0.2941 & ---     & ---      & 0.2059 & \textbf{сгорел} \\
        6 & 0.3529 & 0.5000 & 0.6000 & 0.1471 & обслужен \\
        7 & 0.4118 & ---     & ---      & 0.1882 & \textbf{сгорел} \\
        8 & 0.4706 & 0.6000 & 0.7000 & 0.1294 & обслужен \\
        \hline
    \end{tabular}
\end{table}

Как видно из таблицы~\ref{tab:bit_trace}, биты 0--4 успешно обслуживаются, однако
уже пятый бит (прибывший в $t = 0.2941$~с) должен был бы ожидать до $t = 0.5$~с,
что составляет $0.2059$~с $> t_{\text{lim}}$, поэтому он отбрасывается по тайм-ауту.
Сервер продолжает работу, и в момент $t = 0.5$~с начинает обслуживать бит~6.
Бит~7 также не дожидается освобождения сервера и сгорает в $t \approx 0.5818$~с.
Бит~8 успевает попасть на обслуживание в $t = 0.6$~с.

Дальнейший анализ показывает, что в системе устанавливается периодический режим:
сервер всё время занят, а из каждых 17 поступивших битов ровно 7 отбрасываются,
что обеспечивает среднюю пропускную способность 10 бит/с и интенсивность потерь
7 бит/с. Средняя длина очереди (в штуках) колеблется около значения
$\lambda t_{\text{lim}} \approx 2.89$ бит, однако в дискретной модели это число
реализуется как среднее по времени, а не непрерывный уровень.

\subsection*{Реализация в Simulink}
Для корректного воспроизведения дискретно-событийной природы процесса использована
библиотека SimEvents, позволяющая оперировать отдельными заявками. Основные блоки
и их назначение:

\begin{enumerate}
    \item \textbf{Time-Based Entity Generator} — генерирует биты с постоянной
          интенсивностью $\lambda = 17$~с$^{-1}$.
    \item \textbf{FIFO Queue} — буфер с дисциплиной «первым пришёл — первым обслужен».
    \item \textbf{Single Server} — обслуживает заявки с фиксированным временем
          $1/\mu = 0.1$~с.
    \item \textbf{Event-Based Timeout} (настраиваемый блок) — удаляет заявку из
          очереди, если время её пребывания в буфере превысило $t_{\text{lim}}$.
    \item \textbf{Entity Sink} — поглощает обслуженные биты.
    \item \textbf{Scope, To Workspace} — регистрируют текущую длину очереди и
          накопленное число отброшенных заявок.
\end{enumerate}

Такой подход гарантирует точное соблюдение логики тайм-аута на уровне отдельных
битов, что невозможно получить с помощью непрерывных интеграторов.

\subsection*{Схема модели}
На рис.~\ref{fig:sheme2} представлена блок-диаграмма модели в среде Simulink/SimEvents.

\begin{figure}[h!]
    \centering
    \includegraphics[width=\textwidth]{shema2finel.png}   % <-- скриншот схемы Part 2
    \caption{Схема модели с тайм-аутом в SimEvents.}
    \label{fig:sheme2}
\end{figure}

\subsection*{Результаты моделирования}
На рис.~\ref{fig:queue2} показано состояние буфера $q(t)$ (число битов в очереди,
без учёта обслуживаемого). График имеет ступенчатый характер; максимальная
наблюдаемая длина не превышает 3 бит. Среднее значение за длительный интервал
составляет примерно $2.89$ бит, что согласуется с теоретической оценкой
$\lambda t_{\text{lim}}$.

На рис.~\ref{fig:drop2} приведено накопленное количество отброшенных заявок.
После короткого переходного процесса (около $t = 0.3$~с) кривая практически
линейна; её наклон равен $7$ бит/с, то есть в точности разности
$\lambda - \mu$.

\begin{figure}[h!]
    \centering
    \includegraphics[width=0.8\textwidth]{graph2finelfinel2.png}   % <-- график q(t)
    \caption{Состояние буфера $q(t)$ при тайм-ауте 0.17 с (дискретно-событийная модель).}
    \label{fig:queue2}
\end{figure}

\begin{figure}[h!]
    \centering
    \includegraphics[width=0.8\textwidth]{graph2finelfinel1.png}    % <-- график накопленных отбросов
    \caption{Накопленное число отброшенных заявок.}
    \label{fig:drop2}
\end{figure}

\subsection*{Анализ}
Результаты полностью подтверждают аналитический расчёт:
\begin{itemize}
    \item Первое отбрасывание происходит для бита №~5, как показано в
          таблице~\ref{tab:bit_trace} и видно на графике отбрасываний
          (начало роста кривой).
    \item Установившаяся средняя длина очереди $\approx 2.89$ бит совпадает
          с величиной $\lambda \cdot t_{\text{lim}}$, вытекающей из закона
          Литтла для системы с ограниченным ожиданием.
    \item Наклон кривой отброшенных заявок на линейном участке строго
          соответствует разности $\lambda - \mu = 7$ бит/с, что говорит
          о правильной работе механизма тайм-аута: весь избыток входного
          потока удаляется, а сервер загружен на 100\,\%.
\end{itemize}

\subsection*{Вывод по части 2}
Применение дискретно-событийного подхода (SimEvents) позволило точно
смоделировать поведение системы с индивидуальными тайм-аутами для каждого
бита. Аналитическая трассировка первых заявок выявила момент первого
отбрасывания и причину — превышение допустимого времени ожидания из-за
занятости сервера. Модель продемонстрировала устойчивый режим работы:
очередь стабилизируется около теоретического значения $2.89$ бит,
а интенсивность потерь составляет $7$ бит/с. Полученные результаты
полностью согласуются с теоретическими предсказаниями и подтверждают
корректность выбранной реализации.

% =====================================================================
%  Часть 3 — Непрерывно-детерминированная модель с переменными
%  интенсивностями и ограниченным буфером (Раздел 1, Вариант 1)
%
%  Вставьте этот фрагмент в основной документ ПЕРЕД \end{document}.
%  Перед компиляцией поместите в ту же папку два изображения:
%      scheme3.png  — скриншот схемы Simulink (lab1_variant1)
%      graph3.png   — скриншот окна Scope_Buffer после симуляции
% =====================================================================

\newpage
\section*{Часть 3. Непрерывно-детерминированная модель с переменными интенсивностями и ограниченным буфером}

\subsection*{Цель}
Исследовать динамику изменения объёма данных в буфере узла-приёмника при переменных
во времени скоростях поступления и обработки и при ограниченном размере буфера.
В отличие от Части~1, где интенсивности постоянны, и Части~2, где введён индивидуальный
тайм-аут, здесь основной механизм потерь --- \emph{переполнение буфера ограниченной длины}.

\subsection*{Параметры модели}
В качестве варианта взят порядковый номер $N=17$. Начальные интенсивности $\lambda_{0}$
и $\mu_{0}$ заданы как случайные числа, полученные функцией \verb|rand(1, group_size)|
с фиксированным зерном генератора \verb|rng(N)| для воспроизводимости. Максимальный
объём буфера принят равным порядковому номеру: $K = N = 17$~бит. Параметры сведены
в таблицу~\ref{tab:params3}.

\begin{table}[h!]
    \centering
    \caption{Параметры модели с переменными интенсивностями (Часть~3)}
    \label{tab:params3}
    \begin{tabular}{|l|c|l|}
        \hline
        Параметр & Обозначение & Значение \\
        \hline
        Порядковый номер в группе       & $N$              & 17 \\
        Размер группы                   & $group\_size$    & 20 \\
        Начальная интенсивность входа   & $\lambda_{0}$    & $\approx 0{,}5018$ бит/с \\
        Начальная интенсивность обработки & $\mu_{0}$      & $\approx 0{,}2424$ бит/с \\
        Максимальный объём буфера       & $K$              & 17 бит \\
        Начальное состояние буфера      & $B(0)$           & 0 бит \\
        Шаг по времени                  & $\Delta t$       & 1 с \\
        Время моделирования             & $T$              & 30 с \\
        \hline
    \end{tabular}
\end{table}

\subsection*{Математическая модель}
Скорость поступления изменяется по \emph{линейному} закону, скорость обработки --- по
\emph{логарифмическому} (с округлением до ближайшего целого для сохранения неделимости
бита):
\[
\lambda(t) = \lambda_{0} + t, \qquad
\mu(t) = \operatorname{round}\!\bigl(\log_{2}\!\max(t,1) + \mu_{0}\bigr).
\]
Защита \(\max(t,1)\) необходима, чтобы избежать \(\ln(0)=-\infty\) в первый момент
времени. В модели она реализована блоком \texttt{Saturation} с нижним пределом~1.

Динамика буфера описывается дифференциальным уравнением первого порядка с двумя
кусочно-линейными ограничениями (снизу нулём, сверху размером буфера~$K$):
\[
\frac{dB(t)}{dt} = \lambda(t) - \mu(t), \qquad
0 \leqslant B(t) \leqslant K, \qquad B(0)=0.
\]
При $B(t) > K$ избыточные биты теряются (отбрасывание заявок при переполнении).
При $\lambda(t)<\mu(t)$ и пустом буфере $B(t)=0$ выход «обнуляется» (нельзя
обработать несуществующие биты).

Поскольку шаг $\Delta t = 1$, переход к дискретной форме (схема Эйлера вперёд)
даёт рекуррентное соотношение
\[
B(k+1) = \min\!\bigl(K,\, \max\!\bigl(0,\, B(k) + (\lambda(k) - \mu(k))\cdot \Delta t\bigr)\bigr).
\]

\subsection*{Реализация в Simulink}
Модель построена в соответствии с методикой раздела «Описание и порядок выполнения работы»
(пункты~0--5). Использованы только базовые блоки библиотеки Simulink (без SimEvents):

\begin{enumerate}
    \item \textbf{Digital Clock (Clock)} --- источник времени с шагом $\Delta t = 1$~с;
          \textbf{Rounding Function (Round\_t)} с операцией \texttt{floor} --- округление
          выхода часов до целого, в соответствии с рекомендацией методички.
    \item \textbf{Saturation (Sat\_t)} --- ограничение времени снизу значением~1, что
          исключает $\ln(0)=-\infty$ перед логарифмированием.
    \item \textbf{Constant (lambda0, mu0, two, zero, Kmax)} --- константы для
          $\lambda_{0}$, $\mu_{0}$, основания логарифма~2, нижнего ограничения буфера
          (0~бит) и его максимума ($K$~бит).
    \item \textbf{Sum (Sum\_lambda)} c порядком знаков \texttt{++} реализует
          $\lambda(t)=\lambda_{0}+t$.
    \item Блоки \textbf{Math Function (Log\_t, Log\_2)} с функцией \texttt{log}
          и блок \textbf{Divide} (\texttt{*/}) реализуют переход к основанию~2 по формуле
          $\log_{2}(t)=\ln(t)/\ln(2)$.
    \item \textbf{Sum (Sum\_mu)} прибавляет $\mu_{0}$, после чего
          \textbf{Rounding Function (Round\_mu)} c операцией \texttt{round} даёт
          целочисленное значение $\mu(t)$.
    \item \textbf{Sum (Diff)} с порядком знаков \texttt{+-} вычисляет разность
          $\lambda(t)-\mu(t)$.
    \item \textbf{Discrete-Time Integrator (Buffer)} с шагом $\Delta t=1$ и начальным
          условием $B_{0}=0$ интегрирует разность, реализуя дифференциальное уравнение
          в дискретной форме (Forward Euler, $K\,T_{s}/(z-1)$).
    \item \textbf{Switch (Sw\_nonneg)} с критерием $u_{2}\geqslant 0$ обнуляет
          отрицательные значения буфера (защита от «отрицательных бит»).
    \item \textbf{Switch (Sw\_cap)} с критерием $u_{2}>K$ ограничивает буфер сверху
          значением $K=17$~бит (переполнение --- отбрасывание).
    \item \textbf{Scope (Scope\_Buffer)} и \textbf{To Workspace (B\_log)} ---
          регистрация и экспорт сигнала $B(t)$, подключены \emph{после} обоих
          ограничений, как требует методика.
\end{enumerate}

Все параметры блоков заданы не числами, а именами переменных (\texttt{lambda0},
\texttt{mu0}, \texttt{K}, \texttt{B0}, \texttt{dt}), которые экспортируются в
\texttt{base workspace} из скрипта-сборщика. Это позволяет менять вариант одной
правкой в одном месте, без редактирования полей блоков.

\subsection*{Схема модели}
На рис.~\ref{fig:sheme3} представлена итоговая блок-схема модели Simulink
(\verb|lab1_variant1.slx|).

\begin{figure}[h!]
    \centering
    \includegraphics[width=\textwidth]{shema33333finel.png}      % <-- скриншот схемы Part 3
    \caption{Схема модели Раздела~1, Вариант~1 в среде Simulink.}
    \label{fig:sheme3}
\end{figure}

\subsection*{Результаты моделирования}
На рис.~\ref{fig:queue3} приведён график зависимости состояния буфера $B(t)$ от
времени, полученный с осциллографа \texttt{Scope\_Buffer}.

\begin{figure}[h!]
    \centering
    \includegraphics[width=0.85\textwidth]{graf33333finel.png}   % <-- скриншот Scope_Buffer
    \caption{Состояние буфера $B(t)$ при переменных $\lambda(t)$ и $\mu(t)$.}
    \label{fig:queue3}
\end{figure}

\subsubsection*{Численная трассировка первых шагов}
В таблице~\ref{tab:trace3} приведено пошаговое решение разностного уравнения,
поясняющее ступенчатый характер графика на рис.~\ref{fig:queue3}.

\begin{table}[h!]
    \centering
    \caption{Трассировка первых шагов дискретной модели}
    \label{tab:trace3}
    \begin{tabular}{|c|c|c|c|c|}
        \hline
        $k$ (c) & $\lambda(k)$ & $\mu(k)$ & $\Delta B$ & $B(k{+}1)$ \\
        \hline
        0 & 0{,}50 & 0 & $+0{,}50$ & 0{,}50 \\
        1 & 1{,}50 & 0 & $+1{,}50$ & 2{,}00 \\
        2 & 2{,}50 & 1 & $+1{,}50$ & 3{,}50 \\
        3 & 3{,}50 & 2 & $+1{,}50$ & 5{,}00 \\
        4 & 4{,}50 & 2 & $+2{,}50$ & 7{,}50 \\
        5 & 5{,}50 & 3 & $+2{,}50$ & 10{,}00 \\
        6 & 6{,}50 & 3 & $+3{,}50$ & 13{,}50 \\
        7 & 7{,}50 & 3 & $+4{,}50$ & 18{,}00 $\rightarrow$ \textbf{17 (перепол.)} \\
        8 & 8{,}50 & 3 & --- & 17 \\
        \hline
    \end{tabular}
\end{table}

\subsection*{Анализ}
Из графика и трассировки видны следующие закономерности:

\begin{itemize}
    \item \textbf{Начальный участок} $t\in[0;\,7]$~c. Буфер монотонно растёт ступеньками
          переменной высоты, так как $\Delta B(k)=\lambda(k)-\mu(k)$ увеличивается:
          линейный $\lambda(t)$ растёт быстрее, чем логарифмический $\mu(t)$.
    \item \textbf{Первое переполнение} наступает на шаге $k=8{-}9$~c, когда накопленный
          интеграл превышает $K=17$~бит. Дальнейшие избыточные биты теряются
          (срабатывает блок \texttt{Sw\_cap}).
    \item \textbf{Установившийся режим} $t \gtrsim 9$~c. Буфер удерживается на верхней
          границе $K=17$, потому что разность $\lambda(t)-\mu(t)$ остаётся
          положительной до конца моделирования и интегратор постоянно «прижат» к крышке.
    \item \textbf{Опустошение} ($B=0$) не наблюдается: $\mu(t)$ растёт лишь
          логарифмически и нигде на $[0;\,T]$ не превышает $\lambda(t)$.
\end{itemize}

\subsubsection*{Условия переполнения и опустошения}
Решение поставленной задачи «определить, при каких значениях
$(\lambda,\,\mu)$ буфер переполняется или опустошается» формулируется так:

\begin{itemize}
    \item Буфер \textbf{переполняется}, если
    \[
        \int_{0}^{t}\bigl(\lambda(\tau)-\mu(\tau)\bigr)\,d\tau > K
        \quad \Leftrightarrow \quad
        \lambda(t) > \mu(t) \text{ на достаточно длинном отрезке.}
    \]
    Для Варианта~1 это условие выполняется всегда, начиная с $t \approx 9$~с, т.\,к.
    линейный рост $\lambda(t)$ обгоняет логарифмический рост $\mu(t)$.
    \item Буфер \textbf{опустошается}, если
    \[
        \mu(t) \geqslant \lambda(t) \text{ и } B(t)=0.
    \]
    В рамках Варианта~1 это условие не достигается ни в одной точке промежутка
    моделирования.
\end{itemize}

\subsection*{Вывод по части 3}
Непрерывно-детерминированный подход с использованием классической библиотеки Simulink
(без SimEvents) позволил адекватно описать ИС с переменными во времени интенсивностями
и ограниченным буфером. Модель построена строго в соответствии с методикой:
дифференциальное уравнение реализовано блоком \texttt{Discrete-Time Integrator},
закон $\mu(t)=\operatorname{round}(\log_{2}t+\mu_{0})$ --- последовательностью
\texttt{Math Function~(log) $\rightarrow$ Divide $\rightarrow$ Round} по формуле
приведения логарифма, ограничения сверху и снизу --- двумя блоками \texttt{Switch}.
Результаты моделирования совпадают с численной трассировкой: буфер быстро
заполняется и фиксируется на максимуме $K=17$~бит уже к $t\approx 9$~с,
что подтверждает корректность модели и иллюстрирует асимметрию характеристик
поступления и обработки в системе.


% ======================================================
%  Раздел 2. ДИСКРЕТНО-ДЕТЕРМИНИРОВАННЫЙ ПОДХОД
%  (вставьте этот фрагмент в основной документ перед \end{document})
% ======================================================

\newpage
\begin{center}
    {\Large \textbf{Раздел 2: ДИСКРЕТНО-ДЕТЕРМИНИРОВАННЫЙ ПОДХОД}}
\end{center}
\vspace{10mm}

%----------------------------------------------------------
\subsection*{Задание 1. Описание состояний и условий переходов}

Рассматриваемая система представляет собой одноканальную систему массового обслуживания типа \(M/D/1\) с бесконечным буфером и ограничением на время ожидания заявки в очереди.

Состояние системы определяется величиной \(N(t)\) — общим числом бит в системе в момент времени \(t\). В это число входят биты, находящиеся в очереди, а также бит, который в данный момент обслуживается сервером.

Возможны следующие состояния системы:
\begin{itemize}
    \item \(S_0\) — система свободна, буфер пуст, сервер не занят;
    \item \(S_1\) — один бит находится на обслуживании, очередь отсутствует;
    \item \(S_n\) \((n \ge 2)\) — один бит обслуживается сервером, остальные \(n-1\) бит ожидают в очереди FIFO.
\end{itemize}

В системе происходят три типа событий:
\begin{enumerate}
    \item \textbf{Поступление бита} с интенсивностью
    \[
    \lambda = 17 \text{ бит/с}.
    \]
    При поступлении нового бита система переходит из состояния \(S_n\) в состояние \(S_{n+1}\).

    \item \textbf{Завершение обслуживания} с интенсивностью
    \[
    \mu = 10 \text{ бит/с}.
    \]
    После завершения обслуживания один бит покидает систему, поэтому выполняется переход:
    \[
    S_n \rightarrow S_{n-1}.
    \]

    \item \textbf{Таймаут ожидания}.  
    Если время ожидания бита в очереди превышает допустимое значение
    \[
    t_{\text{lim}} = 0.17 \text{ с},
    \]
    то бит удаляется без обслуживания. В этом случае также происходит переход:
    \[
    S_n \rightarrow S_{n-1}.
    \]
\end{enumerate}

Граф состояний представляет собой бесконечную цепочку состояний, где переходы вправо соответствуют поступлению бит, а переходы влево — завершению обслуживания или удалению бита по таймауту (рис.~\ref{fig:state_graph}).

\begin{figure}[h!]
    \centering
    \includegraphics[width=0.8\textwidth]{digram1.png} 
    \caption{Граф состояний системы M/D/1 с таймаутом.}
    \label{fig:state_graph}
\end{figure}

Поскольку \(\lambda > \mu\), система перегружена, поэтому без механизма таймаута длина очереди возрастала бы неограниченно. Ограничение времени ожидания стабилизирует работу системы за счёт удаления «просроченных» бит.

%----------------------------------------------------------
\subsection*{Задание 2. Блок-схема алгоритма функционирования}

Алгоритм работы системы реализован как дискретно-событийная модель с шагом времени
\[
dt = 0.001 \text{ с}.
\]

В начале моделирования инициализируются основные параметры системы: текущее время, очередь, счётчики поступивших и отброшенных бит, а также состояние сервера.

На каждом шаге моделирования выполняются следующие действия:
\begin{enumerate}
    \item \textbf{Поступление новых бит.}  
    Проверяется момент поступления новых заявок. Все поступившие биты добавляются в очередь.

    \item \textbf{Обслуживание бит.}  
    Если сервер свободен и очередь не пуста, первый бит передаётся на обслуживание. Время завершения обслуживания рассчитывается как:
    \[
    t + \frac{1}{\mu}.
    \]

    \item \textbf{Проверка таймаута.}  
    Для бит, находящихся в очереди, вычисляется время ожидания. Если оно превышает значение \(t_{\text{lim}}\), бит удаляется из системы.

    \item \textbf{Сбор статистики.}  
    После каждого шага фиксируются текущая длина очереди и число отброшенных бит.
\end{enumerate}

Моделирование продолжается до достижения заданного времени \(T_{\text{sim}}\), после чего строятся графики изменения длины очереди и количества потерянных бит. Блок-схема описанного алгоритма приведена на рис.~\ref{fig:block_scheme}.

\begin{figure}[h!]
    \centering
     \includegraphics[width=\textwidth]{flowchart22.png}
    \caption{Блок-схема дискретно-событийного моделирования.}
    \label{fig:block_scheme}
\end{figure}

% ============================================================
%  Вставьте этот фрагмент сразу после Задания 2 (перед \end{document})
%  в файл Раздела 2 "ДИСКРЕТНО-ДЕТЕРМИНИРОВАННЫЙ ПОДХОД"
% ============================================================

\subsection*{Часть 3. Поиск оптимального числа процессоров}
% -----------------------------------------------------------

\textbf{Цель.}  
В реальных системах скорость обслуживания может масштабироваться за счёт
параллельной работы нескольких идентичных процессоров. Требуется подобрать
такое число процессоров \(N_m\), которое минимизирует одновременно и среднее
заполнение буфера, и долю времени простоя оборудования. Исследование проводится
с помощью дискретно-детерминированной модели, построенной в Simulink по методике
раздела~2.

\subsubsection*{Математическая модель}

Интенсивность поступления битов изменяется по закону
\[
\lambda(t) = \lambda_0 + \delta \cdot t \cdot (-1)^t,
\]
а интенсивность обработки одного процессора — по закону
\[
\mu(t) = \operatorname{round}\!\bigl(\mu_0 + \log_2 t \cdot (-1)^t\bigr),
\]
где \(\lambda_0 = 17\) бит/с, \(\mu_0 = 10\) бит/с, \(\delta = 1\), \(t\) — целое
число секунд (шаг моделирования \(\Delta t = 1\) с). Для защиты от \(\ln(0)\) в
момент \(t = 0\) используется ограничение \(t \geqslant 1\).

Система содержит \(N_m\) параллельных процессоров, поэтому суммарная
интенсивность обслуживания равна \(N_m \cdot \mu(t)\). Уравнение состояния
буфера (ёмкость \(K = 17\) бит) записывается как
\[
\frac{dq(t)}{dt} = \lambda(t) - N_m\,\mu(t),\qquad q(0)=0,
\]
с ограничениями \(0 \leqslant q(t) \leqslant K\) (переполнение вызывает потерю
битов).

% ============================================================
%  Вставьте этот фрагмент сразу после Задания 2 (перед \end{document})
%  в файл Раздела 2 "ДИСКРЕТНО-ДЕТЕРМИНИРОВАННЫЙ ПОДХОД"
% ============================================================

\subsection*{Часть 3. Поиск оптимального числа процессоров}

\textbf{Цель.}  
В реальных системах скорость обслуживания может масштабироваться за счёт
параллельной работы нескольких идентичных процессоров. Требуется подобрать
такое число процессоров \(N_m\), которое минимизирует одновременно и среднее
заполнение буфера, и долю времени простоя оборудования. Исследование проводится
с помощью дискретно-детерминированной модели, построенной в Simulink по методике
раздела~2.

\subsubsection*{Математическая модель}

Интенсивность поступления битов изменяется по закону
\[
\lambda(t) = \lambda_0 + \delta \cdot t \cdot (-1)^t,
\]
а интенсивность обработки одного процессора — по закону
\[
\mu(t) = \operatorname{round}\!\bigl(\mu_0 + \log_2 t \cdot (-1)^t\bigr),
\]
где \(\lambda_0 = 17\) бит/с, \(\mu_0 = 10\) бит/с, \(\delta = 1\), \(t\) — целое
число секунд (шаг моделирования \(\Delta t = 1\) с). Для защиты от \(\ln(0)\) в
момент \(t = 0\) используется ограничение \(t \geqslant 1\).

Система содержит \(N_m\) параллельных процессоров, поэтому суммарная
интенсивность обслуживания равна \(N_m \cdot \mu(t)\). Уравнение состояния
буфера (ёмкость \(K = 17\) бит) записывается как
\[
\frac{dq(t)}{dt} = \lambda(t) - N_m\,\mu(t),\qquad q(0)=0,
\]
с ограничениями \(0 \leqslant q(t) \leqslant K\) (переполнение вызывает потерю
битов).

\subsubsection*{Реализация в Simulink}

Модель построена из стандартных блоков. Основные компоненты:

\begin{enumerate}
    \item \textbf{Digital Clock} и \textbf{Rounding Function} (floor) —
          формируют дискретное время \(t\) с шагом 1~с.
    \item Блок \(\boldsymbol{(-1)^t}\) — реализован как \(\cos(\pi t)\).
    \item \textbf{Сумматоры, произведения, блоки Math Function} — собирают
          выражения для \(\lambda(t)\) и \(\mu(t)\) (логарифмирование,
          деление для \(\log_2\), умножение на знак, округление).
    \item \textbf{Gain (Nm)} — умножает интенсивность одного процессора на
          варьируемый коэффициент \(N_m\).
    \item \textbf{Discrete-Time Integrator} (шаг 1 с) — решает разностное
          уравнение с внутренней сатурацией \([0, K]\).
    \item Два блока \textbf{Switch} — гарантируют неотрицательность очереди и
          её ограничение сверху значением \(K\).
    \item \textbf{Scope} и блоки \textbf{To Workspace} — регистрируют сигналы
          для анализа.
\end{enumerate}

Схема модели представлена на рис.~\ref{fig:lab2_scheme}.

\begin{figure}[h!]
    \centering
    \includegraphics[width=\textwidth]{shema2fienl.png}
    \caption{Схема дискретно-детерминированной модели для поиска \(N_m\).}
    \label{fig:lab2_scheme}
\end{figure}

\subsubsection*{Методика перебора}

Число процессоров \(N_m\) перебирается в диапазоне от 1 до
\(\lceil \lambda_0/\mu_0 \rceil + 2 = 4\). Для каждого значения запускается
симуляция длительностью \(T = 40\) с. Регистрируются:
\begin{itemize}
    \item средняя длина очереди \(\langle B \rangle\);
    \item максимальная длина очереди \(B_{\max}\);
    \item доля времени переполнения буфера (overflow);
    \item доля времени простоя процессоров (idle).
\end{itemize}

Оптимальное значение \(N_m^*\) выбирается по минимуму взвешенной суммы
нормированных показателей
\[
J(N_m) = 0{,}5 \cdot (\text{overflow\_frac}) + 0{,}5 \cdot (\text{idle\_frac}),
\]
что обеспечивает баланс между потерями от переполнения и неэффективным
использованием оборудования.

\subsubsection*{Результаты моделирования}

Осциллограф \texttt{Scope\_Buffer} позволяет наблюдать динамику очереди
непосредственно во время симуляции. На рис.~\ref{fig:lab2_graph} приведён
график заполнения буфера для оптимального числа процессоров \(N_m^*\),
полученный с экрана осциллографа.

\begin{figure}[h!]
    \centering
    \includegraphics[width=0.85\textwidth]{graph1forfin.png}
    \caption{Состояние буфера \(q(t)\) при оптимальном \(N_m^*\).}
    \label{fig:lab2_graph}
\end{figure}

\subsubsection*{Анализ и вывод}

Экспериментальный перебор показал, что оптимальным является число процессоров
\(N_m^* = \ldots\) (определяется в ходе симуляции). При этом средняя очередь
составляет около \(\ldots\) бит, максимальная не превышает \(K = 17\) бит, а
доли переполнения и простоя сбалансированы. Дискретно-детерминированный подход
позволил корректно учесть нестационарность потоков и подобрать рациональное
число параллельных обслуживающих устройств.
\end{document}
